\documentclass[11pt]{article}
\usepackage{amsmath,amssymb,amsthm,hyperref}
\newtheorem{theorem}{Theorem}
\newtheorem{proposition}[theorem]{Proposition}
\newtheorem{definition}[theorem]{Definition}
\newtheorem{remark}[theorem]{Remark}
\newcommand{\Q}{\mathbb Q}
\newcommand{\Z}{\mathbb Z}
\newcommand{\R}{\mathbb R}
\newcommand{\C}{\mathbb C}
\newcommand{\CH}{\operatorname{CH}}
\title{GST folds: constructive polarization, integrability, and causal defects}
\author{Research continuation; self-reviewed hand mathematics}
\date{2 November 2026}
\begin{document}
\maketitle

\section{Contract and relation to the existing branch}
Start with an arbitrary original smooth projective complex $X$, an integer
$0\leq p\leq\dim X$, and a prescribed rational Hodge class
$\alpha\in H^{2p}(X,\Q)\cap H^{p,p}(X)$.
The unconditional target remains an independently constructed
$Z_\alpha\in\CH^p(X)_\Q$ with canonical Betti class $\alpha$.
The owner's additional thought experiment assumes this conclusion for all
$X$ temporarily. Results using that assumption are labelled conditional;
they cannot prove the assumption.
Initial checkout: \texttt{c6b8180}. Before publication the branch was fetched
and reconciled through \texttt{05a2758}, including its Round 50--51 records
and Round 51 Lean source. No conclusion is based on their titles alone.

General Space has no dimension field in \texttt{GSTGeneralSpace.lean}.
Actual local Krull dimension is an observation in
\texttt{GSTGeneralSpaceSchemeDimensionObservation.lean}. The higher-cell
index in \texttt{GSTGraphV2UnboundedHigherCausality.lean} is a coherence
index, not automatically a variety's dimension. Here we construct actual
fold operators and test realization, rather than identifying these indices.
The nearest prior fold calculations already distinguish marked Hodge-type
drift, loss of purity outside the smooth-projective domain, and failure of
algebraic realization. We retain that distinction.

\section{A polarization can be constructed from an integral rotation}
\begin{theorem}[Averaged polarization]
Let $\Lambda\simeq\Z^{2s}$ and $J\in\operatorname{End}_{\Z}(\Lambda)$
with $J^2=-1$. Choose any positive definite integral symmetric bilinear
form $b$ on $\Lambda$. On $V=\Lambda\otimes\R$ define
\[
 g_J(x,y)=b(x,y)+b(Jx,Jy),\qquad
 E_J(x,y)=g_J(Jx,y).
 \tag{1}
\]
Then $g_J$ is positive definite and $J$-invariant. The integral form
$E_J$ is alternating, $J$-invariant, and
\[
 E_J(x,Jx)=g_J(x,x)>0\quad(x\ne0).                 \tag{2}
\]
Thus the constant complex structure $J$ on $V/\Lambda$ has an explicit
Riemann form and defines a projective complex torus.
\end{theorem}
\begin{proof}
The equality $J^2=-1$ gives
$g_J(Jx,Jy)=b(Jx,Jy)+b(-x,-y)=g_J(x,y)$.
Consequently $g_J(Jx,y)=-g_J(x,Jy)$, proving alternation using symmetry.
Positivity follows from $g_J(x,x)=b(x,x)+b(Jx,Jx)$; the first summand
is positive for $x\ne0$, and the second is nonnegative.
All lattice values are integral. Constant $J$ is integrable, and (2)
is our symplectic positivity convention. Milne uses the opposite sign,
so his Riemann form is $-E_J$, with $(-E_J)(Jx,x)>0$.
The projectivity implication is Milne, \emph{Abelian Varieties}, v2.00,
Chapter I, Theorems 2.8--2.9 (author-hosted version checked 2 November 2026).
\end{proof}

In a lattice basis with $b(x,y)=x^Ty$, the executable formula is
\[
 G=I+J^TJ,\qquad E=J^TG=-GJ,\qquad
 E^T=-E,\quad J^TEJ=E,\quad EJ=G>0.              \tag{3}
\]
This proves \emph{existence of a chosen polarization}; it does not claim
that a preassigned geometric class $\theta$ becomes positive. In that
precise sense it strengthens, rather than contradicts, Round 31's
remark that its lattice skeleton did not supply a polarization.

\begin{theorem}[Exact geometric range of the constant integral fold]
The complex torus just constructed is isomorphic, as an unpolarized
complex torus and hence as an abelian variety, to $E_i^s$, where
$E_i=\C/(\Z+i\Z)$. In particular its rational Hodge algebra is generated
by actual divisor classes. This constructs algebraic representatives
of every rational Hodge class on this folded realization, in every
codimension, without assuming the Hodge conjecture.
\end{theorem}
\begin{proof}
The action $i\cdot v=Jv$ makes $\Lambda$ a finite torsion-free
$\Z[i]$-module: if $zv=0$ and $z\ne0$, then
$z\bar z\,v=0$, so $v=0$. The Gaussian integers are a PID, whence
$\Lambda\simeq\Z[i]^s$. A Gaussian basis is a complex-linear lattice
isomorphism, giving the asserted torus isomorphism.

For completeness the divisor span is explicit. On $E_i^s$ take the $s$
coordinate point divisors $D_j$, and, for each $j<k$, take
\[
 A_{jk}=(z_j-z_k)^*[0]-D_j-D_k,\qquad
 B_{jk}=(z_j-i z_k)^*[0]-D_j-D_k.                 \tag{4}
\]
The notation means pullback along the indicated algebraic homomorphism
to $E_i$. Their $s^2$ real $(1,1)$ matrices span all Hermitian matrices:
diagonal units, real symmetric off-diagonal units, and imaginary
skew off-diagonal units. They are rational integral cohomology classes;
after complexification they span $H^{1,1}$.
Every basis monomial $u_{i_1}\cdots u_{i_p}
\bar u_{j_1}\cdots\bar u_{j_p}$ is, up to its exterior sign, a product
of $p$ elements $u_i\bar u_j$ of $H^{1,1}$.
Thus products of (4) span all of $H^{p,p}$ over $\C$.
The rational kernel of the Hodge charge operator complexifies to
$H^{p,p}$, since $J$ is rational and the zero-charge condition is
defined over $\Q$. Hence exact rational linear elimination expresses
each rational Hodge target in rational divisor products. Their actual
Chow product is the required representative on this torus.
This is also the known elliptic-power theorem recorded by Milne,
\emph{Shimura varieties and moduli}, p. 519, Step 1, not a novelty claim.
\end{proof}

The existing \texttt{GSTCMFiniteGeometricDegreeDescent.lean} already
provides native divisor-word construction and rational coefficient
descent. The argument above supplies its geometric character expansion
\emph{at this actual torus}, not for an arbitrary original $X$.

\section{The universal rotation exists after doubling, but changes realization}
For any abelian group $H$ define
\[
 J_H(x,y)=(-y,x),\qquad (x,y)\in H\oplus H.       \tag{5}
\]
It satisfies $J_H^2=-1$, without an even-rank hypothesis on $H$.
Moreover
\[
 (J_H-1)(y-x,-x-y)=2(x,y).                       \tag{6}
\]
For a rational vector space, $(J_H-1)^{-1}=-(J_H+1)/2$.
Over $\Z^r$, the equation $(J_H-1)(u,v)=(x,y)$ has a solution iff
each coordinate of $x+y$ is even; then
$u=(y-x)/2$, $v=-(x+y)/2$. Its cokernel is $(\Z/2)^r$.
This is a uniform constructive deck on doubled cohomology, not a
construction of a deck compatible with the original geometric Hodge
decomposition. On a connected original $X$, $H^0(X,\Q)=\Q$ cannot
itself admit a rational operator squaring to $-1$.

For the prescribed $\alpha$, choose integral coordinates $a=m\alpha$
on the free Betti lattice. On the actual auxiliary torus $E_i^{r+1}$
the two homomorphisms $f_\pm(z)=z_0\pm\sum a_jz_j$ construct divisors
$D_\pm=f_\pm^*[0]$ with exact Hermitian matrix difference
\[
 (e_0+a)(e_0+a)^*-(e_0-a)(e_0-a)^*
       =2(e_0a^*+ae_0^*).                       \tag{7}
\]
This is a real target-dependent algebraic source in the folded space.
It encodes \emph{every} rational Betti vector, even a non-Hodge one.
No natural Chow return from this auxiliary torus to $X$ is supplied
by (5)--(7); assigning such a return by its desired value would insert
the missing original-X source. Thus source birth on the new torus is
not source birth on the original variety.

There is also a genuine algebraic realization of (5) for \emph{every}
original $X$ without changing weights: take $Y=X\sqcup X$ and the
signed swap correspondence, whose matrix of graph cycles is
$\left(\begin{smallmatrix}0&-\Delta_X\\\Delta_X&0\end{smallmatrix}\right)$.
Its action on both $\CH^p(Y)_\Q$ and $H^{2p}(Y,\Q)$ is (5), and these
actions commute with the canonical cycle map. The inclusion
$\alpha\mapsto(\alpha,0)$ has the actual component projection as return.
Consequently
\[
 (\alpha,0)\text{ is algebraic on }Y
       \quad\Longleftrightarrow\quad
 \alpha\text{ is algebraic on }X.                \tag{7a}
\]
This all-$X$ construction closes the bare operator-existence question,
while retaining the source problem exactly. Here $J$ is a degree-zero
endomorphism of the \emph{existing} Hodge structure, so its eigenvalues
$\pm i$ can occur within a single $(p,p)$ piece. Reinterpreting those
eigenspaces as $(1,0)$ and $(0,1)$ would change the realization and
weight; it is not the original complex geometry of $X\sqcup X$.

\section{A positive fold equation need not be integrable}
On the real four-torus with periodic coordinates $x_1,\ldots,x_4$, set
$f(x_1)=2+\sin(2\pi x_1)>0$ and define
\[
 J_f\partial_1=\partial_2,\quad J_f\partial_2=-\partial_1,
 \quad J_f\partial_3=f\partial_4,\quad
 J_f\partial_4=-f^{-1}\partial_3.                \tag{8}
\]
It obeys $J_f^2=-1$ and is compatible with the closed symplectic form
$\omega=dx_1\wedge dx_2+dx_3\wedge dx_4$. Indeed for a real vector $v$,
\[
 \omega(v,J_fv)=v_1^2+v_2^2+fv_3^2+f^{-1}v_4^2>0
       \quad(v\ne0).                           \tag{9}
\]
For the convention
$N_J(U,V)=[JU,JV]-J[JU,V]-J[U,JV]-[U,V]$, direct brackets give
\[
 N_{J_f}(\partial_1,\partial_3)
       ={f'\over f}\partial_3.                 \tag{10}
\]
At $x_1=0$ the coefficient is $\pi\ne0$.
Thus even the fold square equation plus positive symplectic geometry
does not force integrability, let alone a projective Hodge structure.
More precisely, within this entire ansatz on a connected periodic
$x_1$-circle, integrability holds if and only if $f$ is constant.
Necessity follows from (10) and $f>0$; sufficiency follows because
constant coefficients give zero brackets. For constant $f$, the complex
coordinates are $x_1+i x_2$ and $x_3+i x_4/f$, giving the actual
projective torus $E_i\times E_{i/f}$. An external deformation parameter
is not a spatial coordinate: this statement does not forbid varying
constant $f$ between different integrable worlds.
This is an actual compact smooth almost-Kaehler realization; it is not
a complex-projective Hodge-conjecture counterexample. For every original
$X$, the product almost-complex structure on $X\times T^4$ retains this
nonzero component. Additional constant blocks give every larger
stabilized dimension without curing it.

\section{Causal defect transport, including cancellation}
Let $J,K,L$ be complex-structure operators on successive real vector
spaces, and let $T,S$ be successive real linear maps. Define
\[
 \mathcal D_{J,K}(T)=KT-TJ.
\]
Then
\[
 \mathcal D_{J,L}(ST)
 =\mathcal D_{K,L}(S)T+S\mathcal D_{J,K}(T).      \tag{11}
\]
Its iterated form for a causal word $T_n\cdots T_1$ is
\[
 J_nT_n\cdots T_1-T_n\cdots T_1J_0
 =\sum_{j=1}^n T_n\cdots T_{j+1}
       (J_jT_j-T_jJ_{j-1})T_{j-1}\cdots T_1.     \tag{12}
\]
The empty word has zero defect. Vanishing per edge implies vanishing
for every finite path. Conversely endpoint vanishing does not force
per-edge vanishing: for $J=\left(\begin{smallmatrix}0&-1\\1&0\end{smallmatrix}\right)$
and $T=\operatorname{diag}(1,-1)$, $JT-TJ\ne0$ but $T^2=1$.
For weight-one realizations this is exactly the Hodge-morphism
condition. In arbitrary cohomological degree one must use the full
Hodge charge/filtration operator, not identify its Weil operator with
a weight-one $J$. The purely algebraic defect formula still applies.

If the owner assumes universal rational Hodge algebraicity, then
Hodge classes obtained under admissible smooth projective operations
have algebraic representatives by that assumption. A marked class
leaving Hodge type under a real fold does not contradict it. Conversely
the absence of drift in (12) is type preservation, not by itself
surjectivity of the canonical cycle-class map.

\section{A constructed native cycle, beyond the Round 51 shape record}
Given an actual cycle-compatible path realization $P$ and a native seed
$Z_0$, define $Z_{n+1}=P_{\rm native}(\gamma_n)Z_n$.
Naturality, by induction, proves
\[
 \operatorname{cl}(Z_n)=P_{\rm Betti}(\gamma_{n-1})\cdots
       P_{\rm Betti}(\gamma_0)\operatorname{cl}(Z_0).       \tag{13}
\]
The new \texttt{nativeCycleTower} gives the actual expression for $Z_n$;
\texttt{nativeCycleTower\_class} proves (13).
\texttt{rides\_always\_gives} in \texttt{GSTLawStackRound51.lean} previously only paired
two supplied propositions in its explicitly labelled shape record.
The added \texttt{GSTLawStackRound51Actual.rides\_always\_gives\_constructively}
uses (13), with no unrelated transport proposition. The existing native
realization and seed remain explicit inputs, not free source coverage.

Coverage of the address ladder $(2p,p,p)$ does not imply coverage of all
class vectors at those addresses. Already on $\mathbb P^1\times\mathbb P^1$,
the two independent divisor classes have the same address $(2,1,1)$.
Tate stabilization of the first does not produce the second. This is
compatible with Round 51's stated open reachability/realization face;
it prevents reading address-orbit coverage as an original-X source theorem.

\section{Evidence and exact proof boundary}
\texttt{GSTFoldRotation.lean}, \texttt{GSTFoldPolarization.lean},
\texttt{GSTCausalFoldDefect.lean}, and
\texttt{GSTVariableFoldIntegrability.lean} contain algebraic proof drafts.
The torus realization, divisor geometry, and vector-field interpretation
are hand-level results, not instantiated scheme/analytic Lean proofs.
Wolfram and an independent Python implementation test exact matrices,
exterior algebras, integral doors, and the variable-fold brackets.
No finite table proves the universal theorems; proofs are above.
No Lean compilation, full semantic review of the repository, or
unconditional original-X Hodge-conjecture solution is asserted.
\end{document}
